\documentclass[10pt,a4paper]{amsart}
\usepackage[margin=19mm]{geometry}
\usepackage{amsmath,amssymb,mathtools}
\usepackage[scr=rsfs,cal=euler]{mathalfa}
\usepackage{xfrac}
\usepackage[unicode,hidelinks,bookmarks=false]{hyperref}
\newcommand{\ie}{i.e.\ }
\newcommand{\cf}{cf.\ }
\newcommand{\resp}{respectively\ }
\newcommand{\Subsection}{\subsection}
\providecommand{\nobreakdash}{\nobreak\hbox{-}}
\setlength{\emergencystretch}{2em}
\multlinegap0pt
\usepackage{bm}
\usepackage{bbm}
\usepackage{dsfont}
\usepackage{xspace}
\usepackage{cancel}
\usepackage{mathtools}
\usepackage{amsthm}
\makeatletter
\@ifundefined{thm}{\newtheorem{thm}{Thm}}{}
\@ifundefined{lem}{\newtheorem{lem}[thm]{Lemma}}{}
\makeatother
\DeclareMathOperator{\PGL}{PGL}
\newcommand{\Q}{\ensuremath{\mathbb{Q} }}
\title[Finite subgroups of maximal order of the Cremona group over ]{Finite subgroups of maximal order of the Cremona group over the rationals}
\author{Ahmed Abouelsaad}
\date{Source: arXiv:2501.18551v2 (CC BY 4.0)}
\begin{document}
\maketitle

\noindent\emph{Source and licence.} Finite subgroups of maximal order of the Cremona group over the rationals. Version arXiv:2501.18551v2 (CC BY 4.0), distributed under the Creative Commons Attribution 4.0 International licence, as recorded in the arXiv licence metadata for this exact version and shown on the abstract page https://arxiv.org/abs/2501.18551v2. Full rights notice: licenses/arxiv-2501.18551-CC-BY-4.0.txt.\par\smallskip

\noindent\textit{Editorial scope.} Counted unit: the Lem whose source label is \texttt{lemm-29} in the article (source0.tex lines 473--480, source label \texttt{lemm-29}), delivered together with the article's own complete original proof of it (source0.tex lines 481--612, one \texttt{proof} environment of 132 lines). No mathematics is written, repaired or re-derived: statement and proof are the article's own text line for line. Required context retained uncounted: element:O3P2 (lines 447--451). Every definition, display and label of those lines is retained unchanged. Omitted from this package and not redistributed: 1--446, 452--472, 613--916. Nothing inside a retained excerpt is omitted or altered: every retained body is delivered exactly as the article's lines are written, so no omission receipt of any kind accompanies this package. Citations: the retained excerpts carry no citation command at all, so no bibliography entry of the article is transcribed and no reference is redistributed. Reference adaptation: none was needed and none was made; every reference inside the delivered unit resolves inside it, so no unresolved reference or citation is produced. Printed slips kept literal: no printed word, symbol, hypothesis or number of the counted statement or of its proof was altered. Layout and class: the article's own class is \texttt{amsart} with packages including changepage, adjustbox, amsfonts, amsmath, amssymb, amsthm, xy, xcolor, ulem, mathabx, enumitem, hyperref, mathtools, extarrows; the wrapper is the lane's pinned \texttt{amsart} editorial wrapper at 10pt on A4, which restates the theorem-like environments the retained text needs and adds the article's own macro definitions that the retained text needs (\texttt{\textbackslash PGL}, \texttt{\textbackslash Q}), with extra wrapper packages (bm, bbm, dsfont, xspace, cancel, mathtools). The delivered numbering is the unit's own. Assets: no retained span contains a figure environment, a graphics-inclusion command, a PDF-page inclusion, a table or a TikZ picture; no image file is copied into this package, the package declares no asset dependency and no image file is redistributed with it. Licence: version arXiv:2501.18551v2 of this article is distributed under the Creative Commons Attribution 4.0 International licence, as recorded in the arXiv licence metadata for this exact version and shown on the abstract page https://arxiv.org/abs/2501.18551v2. A full rights notice is delivered at licenses/arxiv-2501.18551-CC-BY-4.0.txt. Characters: every retained excerpt is pure ASCII, so the delivered engine, XeLaTeX with its default Latin Modern text font, reports no missing character.
\begin{lem}\label{element:O3P2}
    Let $\alpha\in\PGL_2(\Q)$ be an element of order $3$, then $\alpha$ or $\alpha^{-1}$ is conjugate to $\left(\begin{matrix}
       0&-1\\1&-1 
    \end{matrix}\right)$.
\end{lem}

\begin{lem}\label{lemm-29}
Every element $\alpha$ of order $3$ of $\PGL _3(\Q )$ is conjugate to \[B=\left(\begin{matrix}
       0&0  & \lambda \\
          1&0  &0\\
           0&1 &0
    \end{matrix}\right),\]
    for some $\lambda\in\Q \setminus\{0\}$.
\end{lem}

\begin{proof}
    Let $A\in\mathrm{GL}_3(\Q )$ be an element whose class is $\alpha$ in $\PGL _3(\Q )$ has order $3$, then we have $A^3=\lambda\cdot Id$, where $Id$ is the identity matrix and $\lambda\in\Q \setminus\{0\}$. Let $v\in\Q ^3\setminus\{0\}$ be a nonzero vector.
  \begin{itemize}
        \item [(a)] If $v, Av$ and $A^2v$ are three linearly independent vectors, because $A$ sends $v$ to $A\cdot v,~A\cdot v$ to $A^2\cdot v$ and $A^2\cdot v$ to $A^3\cdot v=\lambda \cdot Id\cdot v=\lambda \cdot v$, we find that $A$ is conjugate to the matrix $B$.
        \item [(b)] If $v, Av$, and $A^2v$ are linearly dependent vectors, then we have two cases:
\end{itemize}
        \begin{itemize}
            \item [(b.1)] In case where $v, Av$ are linearly dependent vectors, then $Av=\mu v$, where $\mu \in \Q \setminus\{0\}$. Replacing $A$ by $A/\mu$,we can assume that $\mu=1$. By a coordinate change, $v=\left(\begin{array}{c}
                 1\\
                 0\\0 
            \end{array}\right)$, hence
\[A=\left(
  \begin{array}{ccc}
\begin{matrix}
    1
\end{matrix}
&
\begin{matrix}
    a_{1} & a_{2}
\end{matrix}\\
\begin{matrix}
    0\\0
\end{matrix}&
\begin{matrix}
    C
\end{matrix}
\end{array}
\right),
\]
where $a_i\in\Q ,$ for all $i=1,2$ and $C$ in $\mathrm{GL}_2(\Q )$. With this choice of $A$,

\[A^3=\left(
  \begin{array}{ccc}
\begin{matrix}
    1
\end{matrix}
&
\begin{matrix}
    u & v
\end{matrix}\\
\begin{matrix}
    0\\0
\end{matrix}&
\begin{matrix}
    C^3
\end{matrix}
\end{array}
\right)=Id,
\]
implies to $u=0,~v=0$ and the class of $C$ in $\PGL _2(\Q )$ is of order $3$, then by Lemma \ref{element:O3P2}, the class of $C$ is conjugate to the matrix $\left(\begin{matrix}
       0&-1\\1&-1 
    \end{matrix}\right)$ or its inverse $\left(\begin{array}{cc}
-1 & 1 
\\
 -1 & 0 
\end{array}\right)$. We may then replace $C$ by a conjugate and assume that $C= \mu\cdot\left(\begin{matrix}
       0&-1\\1&-1 
    \end{matrix}\right)$ or $C=\mu\cdot\left(\begin{matrix}
       0&-1\\1&-1 
    \end{matrix}\right)$ ), where $\mu\in\Q^*$. As $A^3=Id$, we obtain $\mu=1$. Hence 
    
    \begin{itemize}
        \item [(b.1.1)] If we assume that
$A=\left(
  \begin{array}{ccc}
\begin{matrix}
    1
\end{matrix}
&
\begin{matrix}
    a_{1} & a_{2}
\end{matrix}\\
\begin{matrix}
    0\\0
\end{matrix}&
\begin{matrix}
    0&-1\\1&-1
\end{matrix}
\end{array}
\right),
$ then $A^3=Id$ for all $a_{1},a_{2}\in\Q $. Moreover, we have $M^{-1}AM=B$, where 
\[M=\left(\begin{array}{ccc}
a_{1}+2 a_{2}-1 & -2 a_{1}-a_{2}-1 & a_{1}-a_{2}-1 
\\
 -3 & 3 & 0 
\\
 -3 & 0 & 3 
\end{array}
\right).\]
 
\item [(b.1.2)] If we assume that
$A=\left(\begin{array}{ccc}
1 & a_{1} & a_{2} 
\\
 0 & -1 & 1 
\\
 0 & -1 & 0 
\end{array}\right)
,$ then $A^3=Id$ for all $a_{1},a_{2}\in\Q $. Moreover, we have $M^{-1}AM=B$, where $\lambda=1$ and 
\[M=\left(\begin{array}{ccc}
a_{1}^{2}+a_{1} a_{2}+a_{2}^{2}+1 & -a_{1}^{2}-a_{1} a_{2}-a_{2}^{2}+1 & 1 
\\
 -2 a_{1}-a_{2} & a_{1}-a_{2} & a_{1}+2 a_{2} 
\\
 -a_{1}-2 a_{2} & 2 a_{1}+a_{2} & -a_{1}+a_{2}
\end{array}\right).\]

    \end{itemize}
   \item [(b.2)] In the case where $v, Av$ are linearly independent vectors but the three vectors are not linearly independent, up to the change of coordinates, we can assume that the two vectors $v$ and $Av$ are 
   \[\left(\begin{array}{c}
                 1\\
                 0\\0 
            \end{array}\right)~\mathrm{and}~\left(\begin{array}{c}
                 0\\
                 1\\0 
            \end{array}\right),\] and since $A$ sends the first one to the second one, and $A^2\cdot v$ is linearly dependent of $v$ and $A\cdot v$. We obtain 
\[
A=\left(
  \begin{array}{ccc}
0 & a_{1} & a_{2} 
\\
 1 & a_{3} & a_{4} 
\\
 0 & 0 & a_{5} 
\end{array}\right).
\]
where $a_i\in\Q ,$ for all $i\in\{1,..,5\}$. one can see that $a_5$ is an eigenvalue of the matrix $A$, and then we can calculate its eigenvector which is $\left(\begin{array}{c}
                 u\\
                 v\\1 
            \end{array}\right)$ for some $u,v\in\Q $. This reduces this case to the previous studied case.
\end{itemize}
\end{proof}

\end{document}
